\documentclass[a4paper, 11pt]{article}

\usepackage[margin=3cm]{geometry} 
\usepackage{amsthm, amsmath, amssymb}
\usepackage{array}
\usepackage{bussproofs}
\usepackage{booktabs}

\theoremstyle{definition}
\newtheorem{definition}{Definition}[section]
\newtheorem{example}[definition]{Example}
\newtheorem{theorem}[definition]{Theorem}
\newtheorem{lemma}[definition]{Lemma}
\newtheorem{prop}[definition]{Proposition}
\newtheorem{corollary}[definition]{Corollary}

% change spacing
\let\olditemize\itemize
\renewcommand{\itemize}{\olditemize\setlength{\itemsep}{0pt}\setlength{\parsep}{0pt}}
\let\oldenumerate\enumerate
\renewcommand{\enumerate}{\oldenumerate\setlength{\itemsep}{0pt}\setlength{\parsep}{0pt}}

% hilbert proofs
\newcounter{dline}
\newenvironment{deduction}[1][\textit{no hypotheses}]
    {\setcounter{dline}{0}%
     \[\begin{array}{r|l@{\qquad}l}
     & #1 \\ \cline{2-2}}
    {\end{array}\]}
\newcommand{\dl}[2]{\stepcounter{dline}\thedline & #1 & \text{#2} \\}

% tables
\newcounter{factrow}
\newenvironment{comparison}[4][p{0.22\textwidth} p{0.36\textwidth} p{0.36\textwidth}]
    {\setcounter{factrow}{0}%
     \begin{center}
     \begin{tabular}{#1}
     \textsc{#2} & \textsc{#3} & \textsc{#4} \\[4pt]}
    {\end{tabular}
     \end{center}}

% ======================================================================
% WORKSHEET MACROS: you should only need the commands in this cheat sheet.
%
%   SCRIPTS (spaces are application, so just type them as on the sheet)
%     \s{S K K a}             inline script, e.g. \s{(\lam{x} x z) y}
%     \lam{x}                 lambda x.  (use inside scripts)
%
%   RUNS (each \st is one step; the step count is printed for you)
%     \begin{run}
%       \from{K a b c}        the starting script (already filled in)
%       \st{a c}              one step
%       \st{a c}[Int]         ... optionally with a type
%     \end{run}
%     Write \st{...} with no space before the brace.
%
%   CERTIFICATES (numbered automatically)
%     \begin{deduction}[p \to q,\ p]       hypotheses (leave off for none)
%       \dl{p \to q}{hypothesis}
%       \dl{q \to p \to q}{\axI{q, p}}     (I) [q, p]
%       \dl{q}{\MP{1}{2}}                  MP 1, 2
%       \ddl{f x}{q}{\MP{1}{2}}            decorated line: script : statement
%     \end{deduction}
%
%   TYPINGS
%     \begin{typing}
%       \ty{K K}{C \to A \to B \to A}{(app) 1, 2}   prints  n. |- K K : ...
%     \end{typing}
%
%   TREES (write the tree bottom-up in reading order: premises first, then
%   the bar that joins them; bussproofs does the layout)
%     \begin{prooftree}
%       \hyp{A \to B}  \dis{A}{1}     open leaf, then leaf [A]^1
%       \impE{B}                      bar with 2 premises
%       \impI[1]{A \to B}             bar discharging 1  (\impI{..} = none)
%     \end{prooftree}
%     Bars: \impE \impI[n] \andI \andEi \andEii \orIi \orIii \orE{i}{j}
%           \botE \allE \allI \exI \exE{n} \eqI \eqE
%     Anything else: \unary{label}{A} \binary{label}{A} \trinary{label}{A}
%     \goal{A}   placeholder: shows A under a \hyp{\vdots}; delete when done
%
%   GAPS   \fillin{}  prints a blank line; type inside the braces to fill it
% ======================================================================
\usepackage{xcolor}
\definecolor{qgrey}{gray}{0.4}
\makeatletter

% questions: the sheet's text in grey, your answer below it
\newenvironment{q}[1]
    {\subsection*{#1}\begingroup\small\color{qgrey}}
    {\par\endgroup\medskip}
\newenvironment{parts}
    {\begin{enumerate}\renewcommand{\labelenumi}{(\roman{enumi})}}
    {\end{enumerate}}
\newcommand{\fillin}[1]{\if\relax\detokenize{#1}\relax\rule[-2pt]{3em}{0.4pt}\else#1\fi}

% scripts: inside \s{..}, \from, \st, \ty and \ddl a typed space is application
\begingroup\catcode`\ =\active%
\gdef {\ifmmode\ \else\space\fi}%
\endgroup
\newcommand{\appspaces}{\catcode`\ =\active}
\newcommand{\s}{\begingroup\appspaces\s@}
\newcommand{\s@}[1]{\endgroup\ensuremath{#1}}
\newcommand{\lam}[1]{\lambda #1.}

% runs
\newcounter{rstep}
\newenvironment{run}
    {\setcounter{rstep}{0}\[\begin{array}{r@{\quad}l@{\qquad}l}}
    {& \ifnum\value{rstep}>0 \text{\small\itshape (\arabic{rstep}~step\ifnum\value{rstep}=1\else s\fi)}\fi \end{array}\]}
\newcommand{\from}{\begingroup\appspaces\from@}
\newcommand{\from@}[1]{\endgroup\@ifnextchar[{\run@line{}{#1}}{\run@line{}{#1}[]}}
\newcommand{\st}{\begingroup\appspaces\st@}
\newcommand{\st@}[1]{\endgroup\stepcounter{rstep}%
    \@ifnextchar[{\run@line{\rightsquigarrow}{#1}}{\run@line{\rightsquigarrow}{#1}[]}}
\def\run@line#1#2[#3]{#1 & #2 & \if\relax\detokenize{#3}\relax\else : #3\fi \\}

% certificates and typings
\newcommand{\axI}[1]{(I)~$[#1]$}
\newcommand{\axII}[1]{(II)~$[#1]$}
\newcommand{\MP}[2]{MP~#1,~#2}
\newcommand{\ddl}{\begingroup\appspaces\ddl@}
\newcommand{\ddl@}[1]{\endgroup\ddl@@{#1}}
\newcommand{\ddl@@}[3]{\dl{#1 : #2}{#3}}
\newenvironment{typing}
    {\setcounter{dline}{0}\[\begin{array}{r@{.\ \ }l@{\qquad}l}}
    {\end{array}\]}
\newcommand{\ty}{\begingroup\appspaces\ty@}
\newcommand{\ty@}[1]{\endgroup\ty@@{#1}}
\newcommand{\ty@@}[3]{\stepcounter{dline}\thedline & \vdash #1 : #2 & \text{#3} \\}

% trees
\newcommand{\hyp}[1]{\AxiomC{$#1$}}
\newcommand{\dis}[2]{\AxiomC{$[#1]^{#2}$}}
\newcommand{\goal}[1]{\noLine\UnaryInfC{$#1$}}
\newcommand{\rl}[2]{\RightLabel{\scriptsize #1\if\relax\detokenize{#2}\relax\else, #2\fi}}
\newcommand{\unary}[2]{\RightLabel{\scriptsize #1}\UnaryInfC{$#2$}}
\newcommand{\binary}[2]{\RightLabel{\scriptsize #1}\BinaryInfC{$#2$}}
\newcommand{\trinary}[2]{\RightLabel{\scriptsize #1}\TrinaryInfC{$#2$}}
\newcommand{\impE}[1]{\binary{$\to$E}{#1}}
\newcommand{\impI}[2][]{\rl{$\to$I}{#1}\UnaryInfC{$#2$}}
\newcommand{\andI}[1]{\binary{$\land$I}{#1}}
\newcommand{\andEi}[1]{\unary{$\land$E$_1$}{#1}}
\newcommand{\andEii}[1]{\unary{$\land$E$_2$}{#1}}
\newcommand{\orIi}[1]{\unary{$\lor$I$_1$}{#1}}
\newcommand{\orIii}[1]{\unary{$\lor$I$_2$}{#1}}
\newcommand{\orE}[3]{\trinary{$\lor$E, #1, #2}{#3}}
\newcommand{\botE}[1]{\unary{$\bot$E}{#1}}
\newcommand{\allE}[1]{\unary{$\forall$E}{#1}}
\newcommand{\allI}[1]{\unary{$\forall$I}{#1}}
\newcommand{\exI}[1]{\unary{$\exists$I}{#1}}
\newcommand{\exE}[2]{\binary{$\exists$E, #1}{#2}}
\newcommand{\eqI}[1]{\AxiomC{}\unary{$=$I}{#1}}
\newcommand{\eqE}[1]{\binary{$=$E}{#1}}

% names
\newcommand{\nm}[1]{\mathit{#1}}
\newcommand{\str}[1]{\text{\texttt{"#1"}}}
\newcommand{\Int}{\mathsf{Int}}
\newcommand{\String}{\mathsf{String}}
\newcommand{\fst}{\mathsf{fst}}
\newcommand{\snd}{\mathsf{snd}}
\newcommand{\Left}{\mathsf{Left}}
\newcommand{\Right}{\mathsf{Right}}
\newcommand{\absurd}{\mathsf{absurd}}
\newcommand{\caseof}[5]{\mathsf{case}\ #1\ \mathsf{of}\ \Left\ #2 \Rightarrow #3 \mid \Right\ #4 \Rightarrow #5}
\makeatother

\begin{document}

\section*{Problem Sheet 4 --- Part B}

% ---------------------------------------------------------------------
\begin{q}{B1. Reading statements}
For (i)--(iv), list the free variables and say which are sentences. For (v)--(viii), read $R(x, y)$ as ``$x$ trusts $y$'',
and say each in plain English.
\end{q}
\begin{parts}
    \item $P(x) \land (\forall x\, Q(x))$: 
    \item $\forall x\, (P(x) \to \exists y\, R(x, y))$: 
    \item $(\exists y\, R(x, y)) \to \forall x\, R(x, y)$: 
    \item $\forall x\, R(x, y)$: 
    \item $\forall x\, \exists y\, R(x, y)$: 
    \item $\exists y\, \forall x\, R(x, y)$: 
    \item $\exists x\, \forall y\, R(x, y)$: 
    \item $\forall x\, \lnot R(x, x)$: 
\end{parts}

% ---------------------------------------------------------------------
\begin{q}{B2. Substitution}
Write out each, or say why it is not allowed.
(i) $A(c/x)$, where $A$ is $P(x) \land (\forall x\, Q(x))$.
(ii) $A(c/x)$ and $A(z/x)$, where $A$ is $\exists y\, R(x, y)$.
(iii) $A(y/x)$, where $A$ is $\exists y\, R(x, y)$. If you put $y$ in blindly, what does the result say, and what should it
say? How do you put it right?
(iv) In Haskell, \texttt{map :: (a -> b) -> [a] -> [b]} has the type $\forall a\, \forall b\, ((a \to b) \to [a] \to [b])$.
The call \texttt{map show [1, 2]} uses it at \texttt{show :: Int -> String}. Which ($\forall$E) bars does the compiler use,
with which $t$, and what type do they give?
\end{q}
\begin{parts}
    \item $A(c/x) = $ 
    \item $A(c/x) = $ \qquad $A(z/x) = $ 
    \item Blindly: \\ Says: \\ Should say: \\ Fix: 
    \item 
\end{parts}

% ---------------------------------------------------------------------
\begin{q}{B3. Reading trees}
For each tree, say which hypotheses are open, write what it shows as ``$\Gamma \vdash A$'', and check the side condition at
each ($\forall$I) and ($\exists$E) bar.
\end{q}
\begin{parts}
    \item {\color{qgrey}\begin{prooftree}
        \hyp{\forall x\, (P(x) \land Q(x))}
        \allE{P(u) \land Q(u)}
        \andEi{P(u)}
        \allI{\forall x\, P(x)}
        \hyp{\forall x\, (P(x) \land Q(x))}
        \allE{P(u) \land Q(u)}
        \andEii{Q(u)}
        \allI{\forall x\, Q(x)}
        \andI{(\forall x\, P(x)) \land (\forall x\, Q(x))}
    \end{prooftree}}
    Open: \qquad Shows: $ \vdash $

    Side conditions: 
    \item {\color{qgrey}\begin{prooftree}
        \hyp{\exists x\, P(x)}
        \dis{P(u)}{1}
        \orIi{P(u) \lor Q(u)}
        \exI{\exists x\, (P(x) \lor Q(x))}
        \exE{1}{\exists x\, (P(x) \lor Q(x))}
    \end{prooftree}}
    Open: \qquad Shows: $ \vdash $

    Side conditions: 
\end{parts}

% ---------------------------------------------------------------------
\begin{q}{B4. ``$\exists$ over $\lor$''}
Fill in the eight gaps of this tree, which shows $\exists x\, (P(x) \lor Q(x)) \vdash (\exists x\, P(x)) \lor (\exists x\, Q(x))$.
Type inside each \verb|\fillin{}|.
\end{q}
{\small
\begin{prooftree}
    \hyp{\exists x\, (P(x) \lor Q(x))}
    \AxiomC{$[\fillin{}]^{\fillin{}}$}
    \dis{P(u)}{2}
    \exI{\fillin{}}
    \orIi{(\exists x\, P(x)) \lor (\exists x\, Q(x))}
    \AxiomC{$[\fillin{}]^{3}$}
    \unary{\fillin{}}{\exists x\, Q(x)}
    \orIii{(\exists x\, P(x)) \lor (\exists x\, Q(x))}
    \trinary{$\lor$E, \fillin{}, \fillin{}}{(\exists x\, P(x)) \lor (\exists x\, Q(x))}
    \binary{$\exists$E, \fillin{}}{(\exists x\, P(x)) \lor (\exists x\, Q(x))}
\end{prooftree}}

% ---------------------------------------------------------------------
\begin{q}{B5. Draw a tree}
Draw a tree with no open hypotheses for each. (iii) The converse of (ii) has no tree. Find two things and a choice of $P$
and $Q$ for which it is false. Where does an attempt at a tree get stuck?
\end{q}
\begin{parts}
    \item \begin{prooftree}
        \hyp{\vdots}
        \goal{(\forall x\, P(x)) \land (\forall x\, Q(x)) \to \forall x\, (P(x) \land Q(x))}
    \end{prooftree}
    \item \begin{prooftree}
        \hyp{\vdots}
        \goal{(\exists x\, (P(x) \land Q(x))) \to (\exists x\, P(x)) \land (\exists x\, Q(x))}
    \end{prooftree}
    \item Two things and $P$, $Q$: \\ Where it gets stuck: 
\end{parts}

% ---------------------------------------------------------------------
\begin{q}{B6. Swapping quantifiers}
Draw a tree with no open hypotheses for each. Hint: in (ii), open both $\exists$s with ($\exists$E) bars near the bottom, with
two different variables.
\end{q}
\begin{parts}
    \item \begin{prooftree}
        \hyp{\vdots}
        \goal{(\forall x\, \forall y\, R(x, y)) \to \forall y\, \forall x\, R(x, y)}
    \end{prooftree}
    \item \begin{prooftree}
        \hyp{\vdots}
        \goal{(\exists x\, \exists y\, R(x, y)) \to \exists y\, \exists x\, R(x, y)}
    \end{prooftree}
\end{parts}

% ---------------------------------------------------------------------
\begin{q}{B7. The old access policy again}
(pol1) $\forall x\, \forall y\, (\mathit{adm}(x) \to \mathit{read}(x, y))$;\enspace
(pol2) $\forall x\, \forall y\, (\mathit{own}(x, y) \to \mathit{adm}(x))$;\enspace
(fact) $\mathit{own}(\mathit{bob}, \mathit{notes})$, Bob owns his notes.
(i) Draw a tree for $\mathit{read}(\mathit{bob}, \mathit{notes})$ whose open hypotheses are (pol1), (pol2) and
$\mathit{own}(\mathit{bob}, \mathit{notes})$. (ii) Add one bar to get $\exists x\, \mathit{read}(x, \mathit{notes})$.
(iii) Write as a statement: ``nobody may read the log, except administrators'', using a name $\mathit{log}$. Why is
$\forall x\, (\mathit{adm}(x) \to \mathit{read}(x, \mathit{log}))$ not it?
\end{q}
\begin{parts}
    \item \begin{prooftree}
        \hyp{\vdots}
        \goal{\mathit{read}(\mathit{bob}, \mathit{notes})}
    \end{prooftree}
    \item \begin{prooftree}
        \hyp{\vdots}
        \goal{\mathit{read}(\mathit{bob}, \mathit{notes})}
        \exI{\exists x\, \mathit{read}(x, \mathit{notes})}
    \end{prooftree}
    \item Statement: $ $ \\ Why not: 
\end{parts}

% ---------------------------------------------------------------------
\begin{q}{B8. ``Nobody may read the vault''}
Can be written $\lnot \exists x\, P(x)$ or $\forall x\, \lnot P(x)$, with $P(x)$ for $\mathit{read}(x, \mathit{vault})$. Show
that the two say the same thing: draw a tree with no open hypotheses for each. Hint: $\lnot A$ is $A \to \bot$: to prove
it, suppose $A$ and aim for $\bot$.
\end{q}
\begin{parts}
    \item \begin{prooftree}
        \hyp{\vdots}
        \goal{\lnot \exists x\, P(x) \to \forall x\, \lnot P(x)}
    \end{prooftree}
    \item \begin{prooftree}
        \hyp{\vdots}
        \goal{(\forall x\, \lnot P(x)) \to \lnot \exists x\, P(x)}
    \end{prooftree}
\end{parts}

% ---------------------------------------------------------------------
\begin{q}{B9. Equality}
Draw a tree with no open hypotheses for each. Why are these two not rules of their own? Hint: for (i), suppose $u = v$. Use
($=$E) with $A$ the statement $z = u$: then $A(u/z)$ is $u = u$, which ($=$I) gives, and $A(v/z)$ is $v = u$.
\end{q}
\begin{parts}
    \item \begin{prooftree}
        \hyp{\vdots}
        \goal{\forall x\, \forall y\, (x = y \to y = x)}
    \end{prooftree}
    \item \begin{prooftree}
        \hyp{\vdots}
        \goal{\forall x\, \forall y\, \forall z\, (x = y \land y = z \to x = z)}
    \end{prooftree}
\end{parts}
Why not rules of their own: 

% ---------------------------------------------------------------------
\begin{q}{B10. Something exists}
Draw a tree with no open hypotheses for each. What does (i) say about the world, and why might that be a surprise?
\end{q}
\begin{parts}
    \item \begin{prooftree}
        \hyp{\vdots}
        \goal{\exists x\, (x = x)}
    \end{prooftree}
    It says: 
    \item \begin{prooftree}
        \hyp{\vdots}
        \goal{(\exists x\, P(x)) \to \lnot \forall x\, \lnot P(x)}
    \end{prooftree}
\end{parts}

% ---------------------------------------------------------------------
\begin{q}{B11. Checking trees}
Each tree below has exactly one mistake. Find the bar that is wrong, and say which condition it breaks.
\end{q}
\begin{parts}
    \item A tree that claims $\forall x\, \lnot \forall y\, (y = x) \vdash \bot$:
    {\color{qgrey}\begin{prooftree}
        \hyp{\forall x\, \lnot \forall y\, (y = x)}
        \allE{\lnot \forall y\, (y = y)}
        \eqI{y = y}
        \allI{\forall y\, (y = y)}
        \impE{\bot}
    \end{prooftree}}
    Mistake: 
    \item A tree that claims $\exists x\, P(x),\ Q(u) \vdash \exists x\, (P(x) \land Q(x))$:
    {\color{qgrey}\begin{prooftree}
        \hyp{\exists x\, P(x)}
        \dis{P(u)}{1} \hyp{Q(u)}
        \andI{P(u) \land Q(u)}
        \exI{\exists x\, (P(x) \land Q(x))}
        \exE{1}{\exists x\, (P(x) \land Q(x))}
    \end{prooftree}}
    Mistake: 
    \item A tree that claims $\vdash \forall y\, \forall x\, (x = y)$:
    {\color{qgrey}\begin{prooftree}
        \eqI{u = u}
        \allI{\forall x\, (x = u)}
        \allI{\forall y\, \forall x\, (x = y)}
    \end{prooftree}}
    Mistake: 
    \item A tree that claims $\exists x\, R(x, u) \vdash \exists x\, R(x, x)$:
    {\color{qgrey}\begin{prooftree}
        \hyp{\exists x\, R(x, u)}
        \dis{R(u, u)}{1}
        \exI{\exists x\, R(x, x)}
        \exE{1}{\exists x\, R(x, x)}
    \end{prooftree}}
    Mistake: 
\end{parts}

% ---------------------------------------------------------------------
\begin{q}{B12. A quantifier inside an implication}
Draw a tree with no open hypotheses for each, and say each in plain English, with $P(x)$ for ``$x$ has signed'' and $Q(x)$
for ``$x$ may vote''.
\end{q}
\begin{parts}
    \item \begin{prooftree}
        \hyp{\vdots}
        \goal{(\forall x\, (P(x) \to Q(x))) \to (\forall x\, P(x)) \to \forall x\, Q(x)}
    \end{prooftree}
    In English: 
    \item \begin{prooftree}
        \hyp{\vdots}
        \goal{(\forall x\, (P(x) \to Q(x))) \to (\exists x\, P(x)) \to \exists x\, Q(x)}
    \end{prooftree}
    In English: 
    \item \begin{prooftree}
        \hyp{\vdots}
        \goal{(\exists x\, (P(x) \to C)) \to (\forall x\, P(x)) \to C}
    \end{prooftree}
    In English: 
\end{parts}

% ---------------------------------------------------------------------
\begin{q}{B13. No side condition (for the curious)}
Suppose ($\forall$I) had no side condition. (i) From the hypothesis $u = c$, derive $\forall x\, (x = c)$ in one bar. Then
derive, with no open hypotheses, first $u = c \to \forall x\, (x = c)$ and then $\forall x\, (x = c)$. (ii) Conclude $\bot$
from the hypothesis $\lnot \forall x\, (x = c)$, which says only that $c$ is not the only thing there is. (iii) Which bar
does the real rule forbid? What would the Haskell compiler call $u$?
\end{q}
\begin{parts}
    \item \begin{prooftree}
        \hyp{u = c}
        \allI{\forall x\, (x = c)}
    \end{prooftree}
    \begin{prooftree}
        \hyp{\vdots}
        \goal{\forall x\, (x = c)}
    \end{prooftree}
    \item \begin{prooftree}
        \hyp{\vdots}
        \goal{\bot}
    \end{prooftree}
    \item 
\end{parts}

% ---------------------------------------------------------------------
\begin{q}{B14. The converse (for the curious)}
B10(ii) showed $(\exists x\, P(x)) \to \lnot \forall x\, \lnot P(x)$. Try to draw a tree with no open hypotheses for the
converse, $\lnot \forall x\, \lnot P(x) \to \exists x\, P(x)$, and say where every attempt gets stuck. Compare A3(b): what
would a tree for it have to do that the hypothesis cannot?
\end{q}
\begin{prooftree}
    \hyp{\vdots}
    \goal{\lnot \forall x\, \lnot P(x) \to \exists x\, P(x)}
\end{prooftree}
Where it gets stuck: 

\end{document}
